\section{Method}
\label{sec:method}

\subsection{Overview and problem setting}

Each EEG segment is represented as a channel--time tensor
$\mathbf{X}\in\mathbb{R}^{C\times T}$, where $C$ is the number of channels
and $T$ is the number of time samples. A minibatch is passed through the same
folding adapter and visual encoder before task-specific prediction. Dataset-specific
sequence handling follows the released evaluation code.

We use standard binary cross-entropy (BCE) for binary tasks and cross-entropy
(CE) for multiclass tasks; label smoothing and other optimization details are
given in Appendix~\ref{app:recipes}. Our contribution lies in the interface
between $\mathbf{X}$ and a generic 2D backbone.

\subsection{Lossless temporal folding}

Choose an integer fold factor $P$ such that $P\mid T$ and define
$W=T/P$.  The folding operator $F_P$ maps the waveform to a single-channel
image $\mathbf{I}\in\mathbb{R}^{1\times(CP)\times W}$:

\begin{equation}
\label{eq:fold}
  \mathbf{I}_{1,cP+p,w}
  = \mathbf{X}_{c,wP+p},
  \quad
  0\leq c<C,\ 0\leq p<P,\ 0\leq w<W.
\end{equation}

Rows are ordered by electrode and then by within-group temporal phase.  Thus
all $P$ phase rows from the same electrode remain adjacent.  For example, one
channel with $T=16$ and $P=4$ becomes

\begin{center}
\small
$\begin{bmatrix}
x_0&x_4&x_8&x_{12}\\
x_1&x_5&x_9&x_{13}\\
x_2&x_6&x_{10}&x_{14}\\
x_3&x_7&x_{11}&x_{15}
\end{bmatrix}$.
\end{center}

Figure~\ref{fig:folding-overview} illustrates the reindexing on a real 16-channel TUEV segment
and a magnified window from the same recording. The raster is a computational
representation, not a physical image of the scalp. It gives the visual encoder
a taller signal-bearing input and places short within-channel temporal
patterns in local 2D neighborhoods before strided stages reduce resolution.
Additional geometric, seizure, and controlled-waveform visualizations appear
in Appendix Figs.~\ref{fig:side-top-view}, \ref{fig:seizure-folding}, and
\ref{fig:synthetic-folding-patterns}.

Because folding is a bijective reindexing, every sample is preserved exactly
and the original waveform is recoverable by the inverse map. Additional
visualizations are provided in the Appendix.

\subsection{Locality induced by folding}

For two folded samples from the same electrode, the image and original-time
distances are respectively
\begin{equation}
d_{\mathrm{img}}=|p-p'|+|w-w'|,\qquad
d_{\mathrm{time}}=|(w-w')P+(p-p')|.
\end{equation}
A vertical step connects consecutive samples, whereas a horizontal step connects
samples $P$ time points apart. Therefore, a local $k\times k$ convolution
captures temporal offsets of the form $\Delta t=P\Delta w+\Delta p$; for
$P=4$, a $3\times3$ kernel combines the offsets
$\{-5,-4,-3,-1,0,1,3,4,5\}$ around its center. Figure~\ref{fig:folding-overview}
visualizes how this kernel captures both adjacent-sample structure and
$P$-spaced temporal context after folding.

\begin{figure*}[t]
\centering
\includegraphics[width=\textwidth]{folding_overview.pdf}
\caption{Lossless folding of a real TUEV training segment ($P=4$).
(a) Sixteen EEG channels before folding; traces share an amplitude scale and
are vertically offset. The highlighted channel and time window identify the
zoom below; the C3--P3 waveform exhibits the original temporal periodicity.
(b) The same preprocessed samples rearranged from $16\times1000$
to $64\times250$; white separators delimit four phase rows per channel.
(The folded view makes the corresponding phase/visual periodic pattern more
explicit.) (c) Sixteen consecutive samples from channel C3--P3, numbered locally from
0 to 15 and colored by phase. (d) The identical samples arranged into four
phase rows: numbers retain the indices in (c). The gold outline encloses a
$3\times3$ kernel centered on sample 5, covering samples 0--2, 4--6,
and 8--10; the small square marks its center. Heatmaps in (b) and (d) share a color
scale. The local window is selected by signal energy for illustration;
folding itself preserves all preprocessed values exactly.}
\label{fig:folding-overview}
\end{figure*}
\FloatBarrier

The layout has two boundary effects: consecutive samples across a phase wrap
are separated across folded columns, and transitions between electrode blocks
depend on channel order rather than physical proximity. Larger $P$ increases
height and temporal-offset diversity but can make these boundaries less local.

\subsection{Architecture-aligned fold-factor selection}
\label{sec:p-selection}

We define a reproducible geometry-based default from the input shape and
backbone stride. For each divisor $P$ of $T$, let

\begin{equation}
  H(P)=CP, \qquad W(P)=T/P.
\end{equation}

EfficientNet-B0 has total spatial stride 32. To retain at least two
signal-bearing cells along height before feature aggregation, we target at least
64 folded rows and choose the smallest valid factor. A native channel count very
close to 64 is treated as sufficient; specifically, SEED-V's 62 channels are
left unfolded. We select

\begin{equation}
\label{eq:p-geometry-rule}
 P_{\mathrm{geo}}=
 \begin{cases}
  1, & C\geq62,\\
  \min\{P\in\mathbb{N}: P\mid T,\ CP\geq64\}, & C<62.
 \end{cases}
\end{equation}

The target height 64 leaves two signal-bearing cells after stride 32. For
ViT-Small, the same folded tensor is zero-padded to the patch-embedding
canvas and tokenized with its native $16\times16$ patch projection; folding
therefore changes the spatial arrangement presented to the patch embedder but
does not alter samples or introduce interpolation. This is
an architecture-alignment heuristic, not a claim that 64 rows are intrinsically
optimal. It uses only $(C,T)$ and backbone geometry and can therefore be fixed
before training. The smallest valid factor limits the loss of horizontal
resolution; task-specific tuning remains possible on validation data.

\subsection{Vision backbone and task head}

All three visual backbones share the folding interface and use a
dataset-specific feature aggregation followed by a linear classifier:

\begin{center}
\fbox{%
\begin{minipage}{0.91\linewidth}
\centering
\(\begin{gathered}
[B,C,T]
\xrightarrow{\;F_P\;}
[B,1,CP,T/P]
\xrightarrow{\;\text{zero pad}\;}
\text{Vision backbone}
\longrightarrow \mathbf{Z}
\\[6pt]
\mathbf{Z}
\xrightarrow{\;\text{GAP / CLS token / Flatten}\;}
\mathbf{z}
\xrightarrow{\;\text{dropout + linear}\;}
\text{logits}
\end{gathered}\)
\end{minipage}}
\end{center}

We use one-channel timm backbones \citep{wightman2019timm}: ImageNet-pretrained
EfficientNet-B0 and DINOv3-pretrained ConvNeXt-Tiny and ViT-Small. Natural-image
classifiers are replaced with task-specific heads, and encoder and head are
fine-tuned jointly. CNNs use global average pooling and ViT uses the CLS token
by default; dataset-specific flatten--FC readouts are listed in
Appendix~\ref{app:recipes}.
